\section{Isolated logarithmic poles on a fibration}
\label{sec:fiber-point-kernels}

Return to a fibration with compact connected fibers, and let the selecting potential have isolated logarithmic poles on the total space. Every selected section descends to the base, and in large tensor powers it must vanish to a certain order at the images of the poles. This order is computed from integrability in all variables of the total space, so it differs from the order imposed by a pole of the same coefficient on the base by at most the fiber dimension. The jets in between span $O(m^{q-1})$ dimensions on a $q$-dimensional base. This is negligible after normalization by $m^q$, so the kernel measures converge to the same quotient equilibrium measure as for the corresponding poles on the base.

The comparison rests on the following elementary fact: for nested section spaces, the total variation of the difference of their kernel measures is the normalized codimension.

\Needspace{19\baselineskip}
\begin{lemma}\label{lem:kernel-codimension-comparison}\index{Bergman kernel!codimension estimate}
Let $B$ be a compact complex manifold of dimension $q$, let $A\to B$ be a holomorphic line bundle with a smooth Hermitian metric, let $v\in C(B)$, and let $\nu$ be a smooth positive probability volume on $B$. For every integer $m\geq1$ let $W_m\subseteq V_m\subseteq H^0(B,A^m)$ be subspaces, both with the norm $\|\cdot\|_{m,v,\nu}$, and put $d_m=\dim V_m-\dim W_m$. Write $B_{V_m,v}$, $B_{W_m,v}$ for their kernels and $\beta_{V_m,v}^{(q)}$, $\beta_{W_m,v}^{(q)}$ for their kernel measures, defined as in \cref{sec:partial-selected-space} with the space indicated by the subscript. Then
\[
 \begin{gathered}
 B_{V_m,v}-B_{W_m,v}\geq0,\\
 \|\beta_{V_m,v}^{(q)}-\beta_{W_m,v}^{(q)}\|_{\mathrm{TV}}
       =\frac{q!d_m}{m^q}.
 \end{gathered}
\]
The determinant functionals of the two spaces, each normalized to vanish at the zero weight, satisfy
\[
 |\mathcal F_{V_m}^{(q)}(v)-\mathcal F_{W_m}^{(q)}(v)|
       \leq\frac{q!d_m}{m^q}\|v\|_\infty.
\]
These statements include the case of zero-dimensional spaces.
\end{lemma}
\begin{proof}
Extend an orthonormal basis of $W_m$ to an orthonormal basis of $V_m$ for the norm at the weight $v$. The difference of the kernels is the sum of the squared weighted evaluations of the $d_m$ added basis vectors. It is nonnegative, and its integral against $\nu$ is $d_m$, which gives the total variation formula. For the determinant functionals, apply the first-variation formula of \cref{prop:partial-kernel-variation} to each of the two fixed spaces, and integrate the difference along $t\mapsto tv$, $0\leq t\leq1$. By what we have just proved, at each $t$ the difference of the kernel measures is positive with total mass $q!d_m/m^q$, so its integral against $v$ is at most $q!d_m\|v\|_\infty/m^q$ in absolute value. Since both functionals vanish at zero, this proves the last assertion.
\end{proof}

Return to a submersion $\pi\colon X\to B$ and a pulled-back line bundle $L=\pi^*A$ as in \cref{sec:partial-submersion}. Equip $A$ with a smooth metric $h_A$ and put $\theta=c_1(A,h_A)$. Choose a smooth positive probability volume $\nu_B$ on $B$, smooth fiber probabilities as in \cref{sec:partial-submersion}, and the smooth positive volume $\nu_X=\mathcal L\nu_B$ on $X$. Use the pullback metric on $L$, and for $v\in C(B)$ use the weight $v\circ\pi$ in the norm on $X$. For a selecting potential $\varphi\in\QPSH(X)$, we write $\beta_{X,m,v\circ\pi,\varphi}^{(q)}$ and $\mathcal F_{X,m}^{(q)}$ for the kernel measure and the determinant functional on $X$ defined by these data, and we regard the selected spaces, by descent as in \cref{prop:partial-kernel-submersion}, as subspaces
\[
 V_m=H^0\bigl(X,L^m\otimes\mathcal I(m\varphi)\bigr)\subseteq H^0(B,A^m).
\]
The following proposition reduces kernel limits for $\varphi$ to the equilibrium theorem on the base, once the selected spaces are close to those of a potential on the base.

\Needspace{24\baselineskip}
\begin{proposition}\label{prop:descended-kernel-limit}\index{Equilibrium measure!descended selected spaces}
Assume that $B$ is K\"ahler, let $\psi\in\PSH(B,\theta)$ have positive non-pluripolar mass, and suppose that
\[
 W_m=H^0\bigl(B,A^m\otimes\mathcal I(m\psi)\bigr)\subseteq V_m,
 \qquad d_m=\dim V_m-\dim W_m=o(m^q).
\]
Let $E(v)=E_{\theta,B}[\psi]_{\mathcal I}(v)$ for $v\in C(B)$. Then, for every $v\in C(B)$, the total variation distance between $\beta_{X,m,v\circ\pi,\varphi}^{(q)}$ and the lift of the kernel measure of $W_m$ on $B$ is $q!d_m/m^q$, and
\[
 \beta_{X,m,v\circ\pi,\varphi}^{(q)}
      \rightharpoonup\mathcal L\bigl(\mathrm{MA}_\theta(E(v))\bigr),
 \qquad
 \mathcal F_{X,m}^{(q)}(v\circ\pi)\longrightarrow
 \int_0^1\!\left(\int_B
           v\,\mathrm{d}\mathrm{MA}_\theta(E(tv))\right)\mathrm{d}t.
\]
The normalized dimensions $q!N_m/m^q$ converge to the total mass of the limit measure, which is finite and positive. If $n>q$, the kernel measures with the normalization $n!/m^n$ tend to zero in total variation.
\end{proposition}
\begin{proof}
For $s\in V_m$, the pointwise norm of $\pi^*s$ and the weight $v\circ\pi$ are constant on the fibers, and every fiber probability has mass one, so the norm of $\pi^*s$ on $X$ equals the norm of $s$ on $B$. As in \cref{prop:partial-kernel-submersion}, pulling back an orthonormal basis shows that $\beta_{X,m,v\circ\pi,\varphi}^{(q)}$ is the lift of the kernel measure of $V_m$ on $B$, and the equality of the Gram matrices shows that $\mathcal F_{X,m}^{(q)}(v\circ\pi)$ is the determinant functional of $V_m$ at $v$. By \cref{lem:kernel-codimension-comparison}, the kernel measures of $V_m$ and $W_m$ differ by a positive measure of total mass $q!d_m/m^q$, and so do their lifts, since lifting preserves positivity and mass; the same lemma bounds the difference of the determinant functionals by $q!d_m\|v\|_\infty/m^q$. Apply \cref{thm:fujiki-partial-kernel-limit,cor:fujiki-partial-determinant} on $B$, with the identity as K\"ahler modification, to the spaces $W_m$ selected by $\psi$. They give the convergence of the kernel measures and determinant functionals of $W_m$, and the total mass of the limit measure is finite and positive and equals the limit of $q!\dim W_m/m^q$. Lifting preserves weak convergence and $d_m=o(m^q)$, which proves the two limits and the assertion about the dimensions. Since the $q$-normalized masses are bounded, multiplying them by $(n!/q!)m^{q-n}$ proves the last assertion.
\end{proof}

We now let the selecting potential have isolated poles. Write $f=n-q$ for the fiber dimension. Fix distinct points $p_1,\ldots,p_N$ of $X$ and positive real numbers $a_1,\ldots,a_N$. Suppose that $\varphi\in\QPSH(X)$ is locally bounded away from these points and that, in local coordinates $x$ centered at $p_j$,
\begin{equation}\label{eq:isolated-selecting-poles}
 \varphi(x)=a_j\log|x|^2+O(1).
\end{equation}
Write
\[
 Z=\{\pi(p_1),\ldots,\pi(p_N)\},\qquad
 a_b=\max_{\pi(p_j)=b}a_j\quad(b\in Z),
\]
and, for an integer $m\geq1$,
\[
 r_{m,b}=\max\{\lfloor ma_b\rfloor-n+1,0\},\qquad
 s_{m,b}=\max\{\lfloor ma_b\rfloor-q+1,0\}.
\]
Let $\mathfrak m_b$ denote the ideal sheaf of a point $b\in B$; its zeroth power is $\mathcal O_B$.

\Needspace{22\baselineskip}
\begin{proposition}\label{prop:isolated-pole-section-descent}
The selected space of $\varphi$ is
\[
 V_m=H^0\left(B,A^m\otimes
                       \bigcap_{b\in Z}\mathfrak m_b^{r_{m,b}}\right).
\]
The space
\[
 W_m=H^0\left(B,A^m\otimes
                       \bigcap_{b\in Z}\mathfrak m_b^{s_{m,b}}\right)
\]
satisfies $W_m\subseteq V_m$ and
\[
 0\leq\dim V_m-\dim W_m
 \leq\sum_{b\in Z}\sum_{j=r_{m,b}}^{s_{m,b}-1}
                         \binom{j+q-1}{q-1}
 =O(m^{q-1}),
\]
where an empty inner sum is zero. The implied constant depends only on the points and their coefficients, not on $m$.
\end{proposition}
\begin{proof}
\Cref{lem:smooth-center-multiplier-ideal}, applied with a point as center, gives
\[
 \mathcal I(m\varphi)_{p_j}=\mathfrak m_{p_j}^{\max\{\lfloor ma_j\rfloor-n+1,0\}};
\]
away from the points $p_j$ the ideal is the full structure sheaf.

Near $p_j$ choose submersion coordinates $(z,\tau)$ with $\pi(z,\tau)=z$. A pulled-back section has a local coefficient $g(z)$ that does not depend on $\tau$, and its vanishing order at $(0,0)$ equals the vanishing order of $g$ at $z=0$. The conditions imposed by the points of one fiber are therefore vanishing conditions at the same base point, and the strongest of them is vanishing to order $r_{m,b}$. This proves the identification.

For the dimension bound, map each section in $V_m$ to its jets modulo $\mathfrak m_b^{s_{m,b}}$ at the finitely many points $b\in Z$. The kernel of this map is $W_m$, so local trivializations of $A^m$ give an injection
\[
 V_m/W_m\hookrightarrow
       \bigoplus_{b\in Z}\mathfrak m_b^{r_{m,b}}/
                                      \mathfrak m_b^{s_{m,b}}.
\]
In $q$ variables there are $\binom{j+q-1}{q-1}$ monomials of degree $j$, which gives the dimensions of these quotients. Finally $0\leq s_{m,b}-r_{m,b}\leq f$ and $s_{m,b}=O(m)$, which proves the bound, also when $q=1$ or $f=0$.
\end{proof}

\Needspace{16\baselineskip}
\begin{theorem}\label{thm:isolated-pole-kernel-limit}
Assume that $B$ is K\"ahler and that there is $\psi\in\PSH(B,\theta)$ with positive non-pluripolar mass, locally bounded off $Z$, such that
\[
 \psi(z)=a_b\log|z-b|^2+O(1)\quad\hbox{near each }b\in Z
\]
in local coordinates $z$ on $B$. Then the selecting potential $\varphi$ of \eqref{eq:isolated-selecting-poles} satisfies the conclusions of \cref{prop:descended-kernel-limit}, and the total variation distance there is $O(m^{-1})$. In particular,
\[
 \beta_{X,m,v\circ\pi,\varphi}^{(q)}
      \rightharpoonup\mathcal L\bigl(\mathrm{MA}_\theta(E_{\theta,B}[\psi]_{\mathcal I}(v))\bigr)
 \qquad(v\in C(B)).
\]
\end{theorem}
\begin{proof}
By \cref{lem:smooth-center-multiplier-ideal} in dimension $q$, the space of sections of $A^m$ selected by $\psi$ on $B$ is the space $W_m$ of \cref{prop:isolated-pole-section-descent}. That proposition gives $W_m\subseteq V_m$ and $\dim V_m-\dim W_m=O(m^{q-1})$, so \cref{prop:descended-kernel-limit} applies.
\end{proof}

If $\theta>0$, \cref{lem:small-isolated-log-poles}, applied on $B$ with the Hermitian form $\theta$, constructs the required $\psi$ whenever the coefficients $a_b$ are sufficiently small. Off the poles, $\theta+\ddc\psi$ then dominates a positive multiple of $\theta$, so the non-pluripolar mass of $\psi$ is positive; it is finite because $\theta$ is closed and $B$ is K\"ahler. The same lemma, applied on $X$ with a Hermitian form, constructs $\varphi$. If the fibers have positive dimension, such a function $\varphi$ does not descend to $B$, since it is finite at the other points of a fiber containing a pole.

The section formula of \cref{prop:isolated-pole-section-descent} also shows that the kernels cannot recover the singularities of the selecting potential. Moving an isolated pole within its fiber does not change any selected space, and neither does adding further poles in that fiber whose coefficients do not exceed the largest one. For the fixed pullback metric $h=\pi^*h_A$, continuous weight $v\circ\pi$ and lifted volume $\nu_X$, all kernels and determinant functionals are therefore the same in every degree. The multiplier ideals on $X$ are nevertheless different: for large $m$, the ideal $\mathcal I(m\varphi)$ of one potential is a proper ideal at a point where the corresponding ideal of the other is the full local ring. Under the hypotheses of \cref{thm:isolated-pole-kernel-limit}, this loss of information occurs together with a positive limit at the growth rate $m^q$. In the closed theory, by contrast, the limit envelope has the multiplier-ideal type of the selecting potential (\cref{thm:hermitian-partial-i-envelope}); this property is lost when the selecting potential is plurisubharmonic only with respect to a Hermitian form.
